\section{Evaluation}
\label{sec:evaluation}

The evaluation is organized around five questions. The first two questions evaluate the recognizer as a source of intent events. The third question evaluates whether the event source can run on the embedded path with usable timing. The last two questions evaluate whether those events can support the safety-interaction story through a control-chain protocol and a user-facing study. This organization keeps recognition quality, runtime feasibility, and safety behavior as separate claims.

\subsection{RQ1: How does the recognizer compare with baselines?}

RQ1 compares the model-quality anchor, the ESP32 runtime candidate, and baseline recognizers under matched noisy-set conditions. The main metrics are accuracy, macro F1, and emergency precision/recall/F1. Emergency is reported explicitly because it is the safety-critical class, while macro F1 keeps movement and unknown from being hidden by a single aggregate score.

\begin{table}[t]
    \centering
    \scriptsize
    \caption{Offline noisy-set recognition evidence. These results support recognizer quality only; they do not validate board-side semantic accuracy or UAV safety behavior.}
    \label{tab:offline_recognition_results}
    \begin{tabular}{@{}p{0.28\linewidth}cccp{0.25\linewidth}@{}}
        \toprule
        Model role & Acc. & Macro F1 & Emerg. P/R/F1 & Claim boundary \\
        \midrule
        Model-quality anchor & 0.88 & 0.88 & 0.96 / 0.79 / 0.87 & Offline noisy-set anchor \\
        ESP32 runtime candidate & 0.85 & 0.85 & 0.95 / 0.70 / 0.81 & Deployment candidate, not safety validation \\
        Baseline recognizers & -- & -- & -- & Add after matched split audit \\
        \bottomrule
    \end{tabular}
\end{table}
% Evidence: model-quality anchor row = weeklyresult/weekly_drone_2026w14/preprocess_ext/classification_report_noisy.txt.
% Evidence: ESP32 runtime candidate row = weeklyresult/weekly_drone_2026w19/xiao_rt1s_c32_b256_samearch_ts/classification_report_noisy.txt.
% TODO(results): Add baseline rows only after verifying common split, support count, frontend, and label mapping.

This table should not become a leaderboard by itself. The final baseline comparison should include compact speech baselines and any alternate frontend variants only after their run configurations are comparable to the anchor. The text should report whether the runtime candidate preserves enough emergency behavior to justify deployment experiments, not whether it replaces the model-quality anchor.

\subsection{RQ2: How robust is the recognizer to rotor noise?}

RQ2 measures how recognition changes as the acoustic condition moves from easier speech to stronger rotor-noisy speech. The expected result is a robustness curve, not a single accuracy number. The curve should report accuracy, macro F1, emergency recall, and unknown false accepts across noise conditions. If the noise is mixed offline, the x-axis can be SNR. If the noise comes from recorded rotor playback or operating conditions, the x-axis should name the condition rather than imply a calibrated acoustic level.

Table~\ref{tab:noise_robustness_plan} is the placeholder for the robustness evidence. The paper should avoid claiming general noise robustness unless each point on the curve has a source result file and a clearly described acoustic condition.

\begin{table}[t]
    \centering
    \scriptsize
    \caption{Rotor-noise robustness evidence plan. Rows become results only after matched source files are attached.}
    \label{tab:noise_robustness_plan}
    \begin{tabular}{@{}p{0.24\linewidth}p{0.30\linewidth}p{0.25\linewidth}p{0.12\linewidth}@{}}
        \toprule
        Condition & What changes & Metrics & Status \\
        \midrule
        Anchor noisy set & Current paper anchor condition & Acc., macro F1, emergency P/R/F1 & Existing \\
        SNR sweep & Offline rotor-noise level & Acc., macro F1, emergency recall, unknown false accepts & TODO \\
        Rotor playback & Recorded or bench-played rotor noise & Same metrics plus condition label & TODO \\
        Live collection & Speaker and drone geometry & Per-condition event accuracy and misses & TODO \\
        \bottomrule
    \end{tabular}
\end{table}
% Evidence: Anchor noisy-set result = weeklyresult/weekly_drone_2026w14/preprocess_ext/classification_report_noisy.txt.
% TODO(results): Add source paths for each SNR, rotor-playback, or live-collection row before reporting the row as a result.

This question is also where frontend comparisons belong. Log-mel, MFCC, and any compact embedded representation should be compared under the same split and reporting format. Table~\ref{tab:frontend_ablation_plan} lists the frontend comparison planned for the draft.

\begin{table}[t]
    \centering
    \scriptsize
    \caption{Frontend-ablation plan for comparing acoustic representations under matched noisy-set conditions.}
    \label{tab:frontend_ablation_plan}
    \begin{tabular}{@{}p{0.25\linewidth}p{0.31\linewidth}p{0.34\linewidth}@{}}
        \toprule
        Frontend & Purpose & Required evidence \\
        \midrule
        Log-mel & Current anchor representation & Matched noisy report and run config \\
        MFCC & Compact speech baseline frontend & Same split, model family, and metrics \\
        Compact embedded frontend & Deployment-oriented representation & Runtime-compatible export and noisy report \\
        Normalized/robust frontend & Noise-robustness comparison & Matched condition labels and report paths \\
        \bottomrule
    \end{tabular}
\end{table}
% TODO(results): Add rotor-noise robustness curve with source paths for each condition.
% TODO(results): Add frontend-ablation table with matched run_config.json and classification_report_noisy.txt paths.
% TODO(boundary): Do not describe untested SNRs, new microphones, or live rotor conditions as validated.

\subsection{RQ3: Can the ESP32 path generate events in real time?}

RQ3 evaluates the embedded event path: audio capture, frontend computation, integer recognizer invocation, and event reporting. Figure~\ref{fig:response_time_breakdown} shows the timing decomposition. The key point is that ``real time'' is not one number. The paper reports board-side invoke time, first-window end-to-end latency, steady-state event period, and BLE inter-arrival time separately because they support different claims.

\begin{figure}[t]
    \centering
    \begingroup
    \setlength{\fboxsep}{3pt}
    \setlength{\tabcolsep}{2pt}
    \resizebox{\linewidth}{!}{%
    \begin{tabular}{c c c c c c c}
        \fbox{\parbox[c][0.52cm][c]{1.35cm}{\centering\scriptsize Audio\\capture}}
        & $\rightarrow$ &
        \fbox{\parbox[c][0.52cm][c]{1.35cm}{\centering\scriptsize Feature\\frontend}}
        & $\rightarrow$ &
        \fbox{\parbox[c][0.52cm][c]{1.35cm}{\centering\scriptsize TFLM\\invoke}}
        & $\rightarrow$ &
        \fbox{\parbox[c][0.52cm][c]{1.35cm}{\centering\scriptsize Event\\report}} \\
        \multicolumn{7}{c}{$\downarrow$} \\
        \fbox{\parbox[c][0.52cm][c]{1.35cm}{\centering\scriptsize Bridge\\decision}}
        & $\rightarrow$ &
        \fbox{\parbox[c][0.52cm][c]{1.35cm}{\centering\scriptsize Command\\boundary}}
        & $\rightarrow$ &
        \fbox{\parbox[c][0.52cm][c]{1.35cm}{\centering\scriptsize Ack /\\timeout}}
        & $\rightarrow$ &
        \fbox{\parbox[c][0.52cm][c]{1.35cm}{\centering\scriptsize State\\log}}
    \end{tabular}%
    }
    \endgroup
    \caption{Response-time evaluation schematic. The paper should report timing by stage and separate board-side inference, first-window end-to-end latency, and steady-state event period.}
    \label{fig:response_time_breakdown}
\end{figure}


\begin{table}[t]
    \centering
    \scriptsize
    \caption{Runtime and transport evidence for the ESP32 path. These measurements support deployment feasibility and event throughput only.}
    \label{tab:runtime_evidence}
    \begin{tabular}{@{}p{0.26\linewidth}p{0.30\linewidth}p{0.25\linewidth}p{0.11\linewidth}@{}}
        \toprule
        Evidence source & Key measurement & Supported claim & Boundary \\
        \midrule
        Local CDC 30-run & 30/30 success; invoke p50 627 ms; total p50 1605 ms & Board-side event generation & Not semantic validation \\
        Dual-core full path & 120/120 success; output p95 1005 ms; first window 1606 ms & Steady-state 1 s throughput & No first-window $<1$ s claim \\
        BLE W30 run & 30/30 success; inter-arrival p50 990 ms; p95 1021 ms & BLE event transport & No flight validation \\
        \bottomrule
    \end{tabular}
\end{table}
% Evidence: Local CDC 30-run = weeklyresult/weekly_drone_2026w19/realworld/esp32_rt1s_c32_validation/validation_summary.json and rt1s_c32_local_cdc_fast30_summary.csv.
% Evidence: Dual-core full path = weeklyresult/weekly_drone_2026w19/realworld/esp32_dual_core_pipeline_validation/mode_c_full_rt1s_summary.json.